% zenodo-paper3-option-a-scoped.tex
% Deltas from main-draft.tex; all unchanged sections omitted.
% Apply by:
%   (1) replacing the \title line (currently main-draft.tex line 59)
%   (2) replacing the \begin{abstract}...\end{abstract} block (lines 74-123)
%   (3) inserting the new \section{Open Problems} block immediately
%       before \section{Conclusion} (currently sec:conclusion, around line 1580)
%
% Rest of the paper (intro, related work, preliminaries, ARE construction,
% composition, provenance, evaluation, security, DORA mapping, conclusion,
% references, appendices) stays intact.
%
% Rationale: the current abstract asserts "we prove" for broader claims
% that the paper's own tab:proven-vs-conjectured marks as Conjectured,
% and that Claim 1 (ARE Pre-Conditioning, Informal) backs with the
% remark "A complete security reduction remains open." This scope-down
% aligns abstract and title with the formal status table.

% ─── (1) Title replacement ──────────────────────────────────────────
% OLD: \title{Certified Heterogeneous Entropy with Algebraic Randomness Extraction}
\title{Certified Heterogeneous Entropy over $\mathrm{GF}(p^n)$:\\
  A Proven Subcase with Conjectured Extensions}

% ─── (2) Abstract replacement ──────────────────────────────────────
\begin{abstract}
Composing entropy from heterogeneous sources, quantum hardware,
WiFi Channel State Information (CSI), and operating system random
number generators, is essential for post-quantum cryptographic
systems but lacks formal provenance guarantees.
%
We present the Certified Heterogeneous Entropy (\che) framework
together with Algebraic Randomness Extraction (\are), a new family
of seeded pre-conditioners parameterized by algebraic programs over
bounded number domains with six arithmetic operations, generated
deterministically from SHAKE-256.
%
\textbf{Proven result.}
For the finite-field domain $\mathrm{GF}(p^n)$, we prove that each
multiplicative step with nonzero operands is a bijection on the
multiplicative group and therefore preserves $\log_2(p^n - 1)$ bits
of min-entropy exactly (\Cref{thm:gf-uniform}, \Cref{cor:gf-chain}).
%
The composition protocol XOR-fuses independent sources with the
standard guarantee that the output retains at least the
min-entropy of the strongest uncompromised input, and records each
contribution in a SHA-256 Merkle-tree certificate that binds the
output to its constituent sources.
%
Graceful degradation adjusts the min-entropy estimate downward
when sources fail, with no silent fallback.
%
\textbf{Conjectured and open.}
A complete extraction soundness reduction for mixed-domain \are
programs remains open; the security of the quaternion
($\mathbb{H}$), octonion ($\mathbb{O}$), and $p$-adic ($\Q_p$)
domains is stated as a conjecture, with the concrete adversarial
models left to future work.
%
The two-layer architecture (algebraic fold followed by SHA-256
expansion) is evaluated empirically but is not a proven extractor
outside the GF subcase.
%
\textbf{Evaluation.}
Our implementation processes 6.8\,MB of IBM Quantum entropy
(ibm\_kingston, 156 qubits), a 9\,KB WiFi CSI pool from ESP32-S3
hardware (treated as a qualitative demonstration, not a certified
source at this scale), and \texttt{os.urandom} output.
%
We report an MCV-based min-entropy proxy; the official NIST
SP~800-90B \texttt{ea\_non\_iid} run is identified as required
future work.
%
The provenance certificates provide an audit trail aligned with
DORA Art.~7 requirements for cryptographic key lifecycle
management in EU-regulated financial institutions.
%
\Cref{tab:proven-vs-conjectured} summarizes the formal status of
each claim; \Cref{sec:open-problems} lists the outstanding proofs
and measurements required to promote the conjectured items.
%
All code is open-source.
\end{abstract}

% ─── (3) New section: Open Problems ────────────────────────────────
% Insert immediately before \section{Conclusion} (sec:conclusion).

\section{Open Problems}
\label{sec:open-problems}

This paper proves a single core theorem (\Cref{thm:gf-uniform}) and
identifies several claims that remain conjectured or require
additional measurement before they meet the same bar. We list them
here so that the extractor theory community and implementers have
a concrete roadmap.

\subsection{Theory gaps}
\label{sec:open-problems:theory}

\begin{enumerate}
  \item \textbf{General \are extraction bound (mixed domain).}
    \Cref{claim:are-extraction} asserts that a random \are program
    of $k$ steps over a source of min-entropy $\minentropy(X) \geq
    \alpha n$ produces output $\varepsilon$-close to uniform for
    suitable $k, \alpha, \varepsilon$. The formal reduction is
    open. The main difficulty is bounding min-entropy across
    non-GF steps (addition, subtraction, multiplication,
    division, modulo, power) composed with the zero-avoidance
    perturbation. \Cref{thm:gf-uniform} is a natural starting
    point because GF steps preserve min-entropy exactly; a bound
    of the form ``if a fraction $\rho$ of steps are GF, loss is
    at most $f(1-\rho)$'' would promote the claim.

  \item \textbf{Quaternion ($\mathbb{H}$) adversary model.}
    The paper asserts that non-commutativity doubles the number
    of distinct evaluation orders the adversary must consider.
    A concrete adversary model (query complexity, success
    probability) is not given. Either a formal argument or a
    softening of the claim to ``structural diversity'' is
    required.

  \item \textbf{Octonion ($\mathbb{O}$) Catalan inversion paths.}
    Same status as $\mathbb{H}$: the Catalan-count argument
    ($C_K$ parenthesizations for $K$ multiplications) is
    combinatorial, not cryptographic. A reduction from a
    standard hardness assumption is open.

  \item \textbf{$p$-adic orthogonal mixing.}
    The $\Q_p$ domain claim ``orthogonal mixing'' in the
    extended security section is unproven. The ultrametric
    structure may or may not yield extraction properties
    distinct from $\R$; no formal result is given.

  \item \textbf{27\% fold collision rate.}
    The 9-domain program layer shows a 27\% collision rate in
    the algebraic fold. The paper does not formally tie this to
    the XOR-composed min-entropy floor. A short argument (XOR
    composition absorbs fold collisions up to the weaker source
    bound) would close the gap, but is not in the current
    manuscript.
\end{enumerate}

\subsection{Measurement gaps}
\label{sec:open-problems:measurement}

\begin{enumerate}
  \item \textbf{NIST SP~800-90B \texttt{ea\_non\_iid} run.}
    The discussion reports an MCV estimator proxy instead of the
    official NIST reference implementation. For a paper that
    invokes DORA Art.~7 audit relevance, the official tool must
    be run on all three sources (IBM Quantum 6.8\,MB, CSI, OS)
    and reconciled with the MCV proxy. The tool is freely
    available from NIST; the blocker is installation and
    execution, not access.

  \item \textbf{CSI pool size.}
    The CSI input in this paper is 9\,KB, below any reasonable
    NIST SP~800-90B entropy-assessment floor (the reference
    recommendation is $\geq 10^6$ samples). Paper~2 uses a
    larger BCM4339 Nexmon trace; either reuse that trace or
    explicitly scope CSI as a qualitative demonstration (as
    this revision does in the abstract) rather than a certified
    source.

  \item \textbf{Cross-source independence.}
    The XOR composition proof assumes source independence. The
    paper does not empirically measure mutual information
    between IBM Quantum, CSI, and OS entropy streams. A simple
    pairwise mutual-information test over the collected traces
    would either validate or bound the independence assumption.
\end{enumerate}

\Cref{tab:proven-vs-conjectured} already lists the formal status
of each claim. This section translates that status into an
actionable research program: items in
\Cref{sec:open-problems:theory} require new proofs; items in
\Cref{sec:open-problems:measurement} require instrument runs on
existing or expanded data. Together they constitute the path to
promote the current manuscript from ``proven GF subcase with
conjectured extensions'' to ``proven mixed-domain \are extractor
with certified source characterization.''
